%!TEX root = ../main/main.tex



En esta pregunta usted va a implementar un sistema de cifrado simétrico basado en la redes de Feistel.
Al igual que en el caso de AES, este sistema de cifrado se define para bloques de un cierto largo, y luego se utilizan modos de operación que indican cómo se cifra un mensaje de largo arbitrario usando el cifrado de bloques. Adicionalmente, y al igual que para AES, se debe utilizar una función de padding para que el largo del mensaje a cifrar sea divisible por el largo del bloque.

\subsection*{Definición del sistema de cifrado}
Considere un largo de bloque $2n$ (el largo de un bloque debe ser par), suponga dada una función de hash $h : \{0,1\}^* \to \{0,1\}^n$, y fije un parámetro $r \geq 1$ como el número de rondas de la red de Feistel. Dados estos parámetros, una red de Feistel define un sistema de cifrado simétrico donde $\M = \C = \{0,1\}^{2n}$ y $\K = \{0,1\}^*$ (vale decir, se puede usar una clave de largo arbitrario). Dado un mensaje $m \in \{0,1\}^{2n}$ y una clave $k \in \{0,1\}^*$, la función de cifrado $\Enc_k : \{0,1\}^{2n} \to \{0,1\}^{2n}$ para un bloque se define de la siguiente forma. El mensaje se divide como $m = m_1 \| m_2$, donde $m_1 \in \{0,1\}^n$ y $m_2 \in \{0,1\}^n$, y se define la siguiente secuencia $\{L_i, R_i\}_{i \in \{0, \ldots, r\}}$ de strings: $L_0 = m_1$, $R_0 = m_2$, y para cada $i \in \{0, \ldots, r-1\}$:
\begin{align*}
L_{i+1} \ &= \ R_i,\\
R_{i+1} \ & = \ L_i \oplus h(k_{i+1} \| R_i),
\end{align*}
donde $k_{i+1} = h(k \| i+1)$ (esto es un key schedule, como el que es usado en AES). Finalmente, se tiene que $\Enc_k(m) = L_r \| R_r$.

\subsection*{La clase \texttt{Feistel} a implementar}
En esta pregunta usted debe implementar en Python una clase \texttt{Feistel} que define un sistema de cifrado simétrico basado en la redes de Feistel. Esta clase debe tener los siguientes componentes:
\begin{itemize}
\item Un atributo \texttt{block\_length} que almacena el largo del bloque, el cual debe ser un número par.

\item Un atributo \texttt{number\_rounds} que almacena el número de rondas de la red de Feistel, el cual debe ser mayor o igual que 1.

\item Un atributo \texttt{hash\_function} que almacena la función de hash que se usará en la red de~Feistel.

\item Un atributo \texttt{secret\_key} que corresponde a la clave secreta que se usará en el cifrado y descifrado de mensajes.

\item Un constructor que recibe el largo del bloque, el número de rondas, la función hash que se utilizará en la red de Feistel, y la clave secreta.

\item Una función \texttt{enc\_block} que cifra un mensaje \texttt{m} de largo \texttt{block\_length} utilizando la clave \texttt{secret\_key}. \texttt{enc\_block} implementa la función
  $\Enc_{\texttt{secret\_key}}$ descrita anteriormente, y no realiza padding para calcular el cifrado.

\item Una función \texttt{dec\_block} que descifra un mensaje \texttt{c} de largo \texttt{block\_length} utilizando la clave \texttt{secret\_key}. \texttt{dec\_block} implementa la
  inversa de la función $\Enc_{\texttt{secret\_key}}$ descrita anteriormente.

\item Una función \texttt{key\_schedule} que genera la clave para la ronda \texttt{i} de la red de Feistel. Para implementar esta función use la instrucción
\texttt{str(i).encode()} para transformar \texttt{i} en una secuencia de bytes.

\item Una función \texttt{pad} que agrega a un mensaje \texttt{m} un símbolo \texttt{1} seguido de un número arbitrario de símbolos \texttt{0} hasta que el largo del string generado es divisible por \texttt{block\_length}. Si el largo de \texttt{m} ya es divisible por \texttt{block\_length}, entonces se agrega un nuevo bloque a \texttt{m} de la forma \texttt{10...0} de largo \texttt{block\_length}.

\item Una función \texttt{unpad} que elimina los símbolos extra agregados a un mensaje \texttt{m} mediante la función \texttt{pad}.

\item Una función \texttt{enc} que recibe como parámetros un mensaje \texttt{m} de largo arbitrario y un modo de operación \texttt{mode}. Esta función cifra el resultado de aplicar la función \texttt{pad} a \texttt{m} usando la clave \texttt{secret\_key} y considerando el modo de operación \texttt{mode}. Los dos modos de operación posibles son \texttt{ECB} y \texttt{CBC}, los cuales son definidos como en el caso de AES. Si el modo de operación es \texttt{ECB}, entonces \texttt{enc} retorna el mensaje cifrado. Si el modo de operación es \texttt{CBC}, entonces \texttt{enc} genera de manera aleatoria un vector de inicialización \texttt{IV} que utiliza para cifrar, y retorna una tupla con el mensaje cifrado y el vector de inicialización \texttt{IV}.

\item Una función \texttt{dec} que recibe como parámetros un mensaje \texttt{c} de largo arbitrario, un modo de operación \texttt{mode}, y un parámetro opcional \texttt{IV} que representa un vector de inicialización. Esta función descifra el mensaje \texttt{c} usando la clave \texttt{secret\_key} y considerando el modo de operación \texttt{mode}, y luego utiliza la función \texttt{unpad} para eliminar los símbolos extra agregados al cifrar un mensaje. Los dos modos de operación posibles son \texttt{ECB} y \texttt{CBC}. Ambos modos son definidos como en el caso de AES, y para el caso de \texttt{CBC} es necesario incluir el vector de inicialización \texttt{IV} como parámetro.
\end{itemize}
Todos los mensajes y claves en su implementación deben ser de tipo \texttt{bytes}. En particular, cuando se indica que un bloque tiene largo 32, significa que el bloque esta formados por 32 bytes, lo cual equivale a 256 bits. Todas las operaciones deben realizarse sobre bytes. Por ejemplo, la operación $a \oplus b$ se debe aplicar sobre dos bytes $a$ y $b$, lo cual en Python es equivalente a transformar $a$ y $b$ en sus representaciones binarias $a'$ y $b'$, respectivamente, realizar $a' \oplus b'$ bit a bit para obtener un resultado $c$ en binario, y finalmente transformar $c$ a un byte. Además, debe utilizar \texttt{raise} para generar excepciones cuando hay errores en los parámetros utilizados; por ejemplo, si el largo del bloque no es un número par, si el número de rondas es igual a 0 o si se intenta descifrar un mensaje usando \texttt{CBC} sin un vector de inicialización \texttt{IV}. Finalmente, se debe utilizar la librería \texttt{typing} para especificar los tipos de los parámetros y los valores retornados por las funciones.

\subsection*{Detalles de la implementación de la clase \texttt{Feistel}}
En esta tarea usted debe completar la siguiente implementación de la clase \texttt{Feistel}, donde además puede agregar las funciones auxiliares que estime conveniente.

\begin{python}
class Feistel:
    def __init__(
        self,
        block_length: int,
        number_rounds: int,
        hash_function: Callable[[bytes], bytes],
        secret_key: bytes,
    ) -> None:
        pass

    def enc_block(self, m: bytes) -> bytes:
        pass

    def dec_block(self, c: bytes) -> bytes:
        pass

    def key_schedule(self, i: int) -> bytes:
        pass

    def pad(self, m: bytes) -> bytes:
        pass

    def unpad(self, m: bytes) -> bytes:
        pass

    def enc(self, m: bytes, mode: str) -> bytes | Tuple[bytes, bytes]:
        pass

    def dec(self, c: bytes, mode: str, IV: Optional[bytes] = None) -> bytes:
        pass
\end{python}

\subsection*{Verificación de la correctitud de la implementación}
Para probar si su implementación es correcta se va a utilizar la siguiente implementación de la función de hash \texttt{SHA256}:
\begin{python}
from hashlib import sha256

def SHA256(message: bytes) -> bytes:
    alg = sha256()
    alg.update(message)
    return alg.digest()
\end{python}
Considerando esta función, usted puede verificar si su implementación funciona correctamente utilizando pruebas como la siguiente, donde se especifica que los bloques son de 64 bytes (puesto que \texttt{SHA256} produce un hash con 32 bytes), y se indica que la red de Feistel tiene 10 rondas:
\begin{python}
net = Feistel(64, 10, SHA256, b'T.d66w7ydubyw3h-Hsy')

m1 = b'1234567890123456789012345678901234567890123456789012345678901234'
c1 = net.enc_block(m1)
d1 = net.dec_block(c1)

m2 = b'Este es un mensaje de largo arbitrario'
c2 = net.enc(m2, 'ECB')
d2 = net.dec(c2, 'ECB')

m3 = b'Este es un mensaje de largo arbitrario'
(c3, iv) = net.enc(m3, 'CBC')
d3 = net.dec(c3, 'CBC', iv)

print(f'Texto plano      : {m1}')
print(f'Texto cifrado    : {c1}')
print(f'Texto descifrado : {d1}\n')
print(f'Texto plano      : {m2}')
print(f'Texto cifrado    : {c2}')
print(f'Texto descifrado : {d2}\n')
print(f'Texto plano      : {m3}')
print(f'Texto cifrado    : {c3}')
print(f'IV               : {iv}')
print(f'Texto descifrado : {d3}')
\end{python}
El resultado de utilizar este código debería ser el siguiente:
{\small
\begin{verbatim}
Texto plano      : b'1234567890123456789012345678901234567890123456789012345678901234'
Texto cifrado    : b'0\xa6\x9c)\xe4\xf9\x04\xbb\xb7\xebS\x98^\x1c5\x1at1\xd7\xc8\xde
                     \xf4\xac\x8e\x08\x95/\x07\x0f\xac\xc2\xf5\xba\x89\xf3Z\xe1\x12\xf5
                     \x10o\xea\xe2!\xe0\xfe\xfd++\x01\xd7\xb8C\xa9\xc9$\xcb\xbfG\xfdj\xa8uq'
Texto descifrado : b'1234567890123456789012345678901234567890123456789012345678901234'

Texto plano      : b'Este es un mensaje de largo arbitrario'
Texto cifrado    : b'\xf7\x80.~\x01\xb7\xf2\xc4\xd6\xb1I\x125\xae\xe3\xb9\x8dq\xbcP\x14
                     \xf7\x84\xaf\x94\x83\xa9%\x17\xa2?\xeet\xc9=\xb0L+\x08\x04q\x14J\x06
                     \\\x18\xbc\xa9t\xdc\x87\xe0\xffx\xa0s\xb7\x17\xca\x08\xd0_\xde\xdb'
Texto descifrado : b'Este es un mensaje de largo arbitrario'

Texto plano      : b'Este es un mensaje de largo arbitrario'
Texto cifrado    : b'V.K\x04\x92\x84\x16yS\xccI\xf1\x01\xaaX\x1a\x8fW{\x8e=HaJ\x85\x12$\xa8
                     \x93;X\xf5\x98\x8c\xaa8\x82Az\x14/l'\x10\xdf\t{\xcfk\xcd\xba\ng\xed6
                     \xbe\xeb\xa8,\x15\xe3\x9e\xdb@'
IV               : b'K\xe55\x7fng\x9cJ\x8c\xbez\xd5k\xe0\x1e\xe2\x19\x92\xe3u\xc5\x88/#\xc4
                     !G(\x16\xf8w\xc6\x96\x94tS*|+\x16e6\xf0\x1c\x96\x86\xdbZp\xe1\xf3(\xdf
                     \xb6iE[\xfa\x1f\xdd6\xae\x07\xc0'
Texto descifrado : b'Este es un mensaje de largo arbitrario'
\end{verbatim}
}
